\begin{center}
	\large{\textbf{STUDI KLASIFIKASI KETERBELITAN TIGA QUBIT 
MENGGUNAKAN PENGISIAN KONKURENSI}}
\end{center}
%\begin{flushleft}
\begin{minipage}{4cm}
	\begin{align*}
		&\text{Nama} \qquad &&: \text{Tamara Pingki} \\
		&\text{NRP} \qquad &&: \text{6001241001} \\
		&\text{Departemen} \qquad &&: \text{Fisika} \\
		&\text{Pembimbing} \qquad &&: \text{Prof. Agus Purwanto, D.Sc.} \\
	\end{align*}
\end{minipage} \\
\begin{center}
	\large\textbf{ABSTRAK} \\  
\end{center}
Keterikatan multipartit sejati merupakan fenomena fundamental dalam mekanika kuantum yang berperan penting dalam teori dan aplikasi informasi kuantum. Namun, karakterisasi dan pengukuran keterikatan multipartit sejati masih menghadapi tantangan teoretis, terutama pada sistem multipartit dan keadaan campuran. Salah satu pendekatan yang berkembang adalah pendekatan geometris, yang memanfaatkan struktur ruang keadaan kuantum dalam mengukur keterikatan. Penelitian ini mengkaji pengisian konkurensi (\textit{concurrence fill}) sebagai ukuran geometris keterikatan multipartit sejati pada sistem tiga-qubit. Analisis difokuskan pada formulasi matematis pengisian konkurensi serta keterkaitannya dengan tiga kekusutan dan kekusutan parsial, sehingga memungkinkan klasifikasi keterikatan pada keadaan kelas GHZ dan W. Selain itu, penerapan pengisian konkurensi pada keadaan campuran tertentu, khususnya campuran peringkat-2 dari keadaan kelas GHZ dan W, juga dibahas untuk menelaah struktur keterikatan yang muncul. Kajian ini dilakukan secara teoretis dan analitis, tanpa melibatkan simulasi numerik maupun verifikasi eksperimental.

\noindent
\begin{minipage}{1cm}
	\begin{align*}
		\text{\textbf{Kata kunci}}:\hspace{5mm}&\text{Keterbelitan, Tiga Qubit, Pengisian Konkurensi}% \\ &\text{} \\&\text{}
	\end{align*}
\end{minipage}

\newpage
\thispagestyle{empty}
\begin{center}
	\topskip0pt
	\vspace*{\fill}
	\textit{\textbf{"Halaman ini sengaja dikosongkan"}}
	\vspace*{\fill}
\end{center}
\pagebreak